\newvideofile{synchronmaschine}{Synchronmaschine}
\v{
	\begin{frame}{Lernziele}
		\begin{Lernziele}{Synchronmaschine}
			Du kannst
			\begin{itemize}
				\item Aufbau und Hauptkomponenten sowie die Vor- und Nachteile einer Synchronmaschine benennen.
				\item erklären, wie das dreiphasige Statorfeld entsteht und warum Rotor und Drehzahl synchron zur elektrischen Frequenz laufen.
				\item den Polradwinkel Theta und seine Auswirkung auf das Drehmoment beschreiben.
			\end{itemize}
		\end{Lernziele}
		\speech{LernzieleSynchronmaschine}{1}{In diesem Kapietel betrachten wir eine weitere, für viele Zwecke eingesetzte, elektrische Maschine: Die Synchronmaschine. Dabei gehen wir auf den Aufbau und die Vor- und Nachteile dierser elektrischen Maschinen ein. Wir werfen eine Blick auf das Statorfeld, das aus drei Phasen besteht und letztlich den Rotor in Bewegung versetzt. Außerdem thematisieren wir dabei den Zusammenhang aus dem Polradwinkel und dem Drehmoment der Synchronmaschine.}
		\silence{2}
	\end{frame}
}
\section{Synchronmaschine}
\s{Die Synchronmaschine ist eine Drehfeldmaschine, die folglich mit einem dreiphasigen Drehstrom (siehe Modul 7) im Statorkreis betrieben wird. Der Rotor kann entweder permanent (mit Permanentmagneten) oder durch einen Gleichstromkreis erregt werden. Dieser Maschinentyp wird in Kraftwerken zur Stromgewinnung eingesetzt, da er aber über einen sehr hohen Wirkungsgrad verfügt, wird er auch in der modernen Antriebstechnik in kleineren Leistungsklassen immer häufiger eingesetzt. Ein wesentlicher Nachteil besteht jedoch darin, dass die Synchronmaschine nur mithilfe einer Leistungselektronik betrieben werden kann, da sie an einem einfachen Drehstromnetz nicht anläuft.}
\vspace{0.5cm}
\v{
	\begin{frame}
		\f{}{height=0.8\textheight}{Aufbau_Synchronmaschine1}{}
	\end{frame}
	\speech{EinfuehrungSynchronmaschine}{1}{Die Synchronmaschine ist eine Drehfeldmaschine, die folglich mit einem dreiphasigen Drehstrom im Statorkreis betrieben wird. Der dreiphasige Drehstrom wrde bereits im Modul 7 behandelt. Der Rotor kann entweder permanent, also mit Permanentmagneten, oder durch einen Gleichstromkreis erregt werden. Dieser Maschinentyp wird in Kraftwerken zur Stromgewinnung eingesetzt, da er aber über einen sehr hohen Wirkungsgrad verfügt, wird er auch in der modernen Antriebstechnik in kleineren Leistungsklassen immer häufiger eingesetzt. Ein wesentlicher Nachteil besteht jedoch darin, dass die Synchronmaschine nur mithilfe einer Leistungselektronik betrieben werden kann, da sie an einem einfachen Drehstromnetz nicht anläuft.}
	\silence{2}
}

\begin{frame}
	\ftx{Synchronmaschine}
	\b{Die Synchronmaschine dreht netzsynchron mit konstanter Drehzahl.\\
	\vspace{0.5cm}}
	\textbf{Vorteile:}
	\begin{itemize}
		\item Robust und wartungsfrei.
		\item Hohe Regeldynamik.
		\item Hohe Leistungsdichte bei geringem Bauvolumen.
		\item Überlastfähigkeit.
	\end{itemize}\pause
	\textbf{Nachteile:}
	\begin{itemize}
		\item Aufwendige Umrichtertechnik.
		\item Aufwendige Regelung.
		\item Gesteuerter Betrieb nur mit großem Aufwand oder gar nicht möglich.
	\end{itemize}
	\speech{SynchronmaschineVorNachteile}{1}{Die Synchronmaschine dreht netzsynchron mit konstanter Drehzahl, daher ihr Name. Sehen wir uns nun zunächst die Vor- und Nachteile einer Synchronmaschine an. Die Vorteile liegen insbesondere in der Robusten Bauweise, weswegn sie sehr wartungsarm ist. Sie hat eine Hohe Regeldynamik, was bedeutet, dass sie schnell und genau auf Last- oder Sollwertänderungen reagiert und Frequenz sowie Drehmoment stabil hält. Sie hat bei einer kleinen Baugröße eine hohe Leisutngsdichte und ist zudem Überlastfähig, was bedeutet, dass sie kurzfristig über ihrem Nennstrom oder ihrer Nennleistung hinaus betrieben werden kann, ohne Schaden zu nehmen, um etwa Spitzenlasten zu decken.}
	\silence{2}
	\speech{SynchronmaschineVorNachteile}{1}{Die Nachteile de Synchronmaischine liegen vor allem in der dafür notwendigen, aufwändeigen Regeltechnik. Der Betrieb erfordert komplexe und teure Leistungselektronik, also Umrichter, um Frequenz und Spannung passend bereitzustellen. Zudem ist die Regelung der Maschine anspruchsvoll, weil Drehzahl, Erregung und Last synchron abgestimmt und dynamisch geregelt werden müssen. Außerdem ist ein gezielt steuerbarer Betrieb, zum Beispiel schnelle Drehzahl- oder Drehmomentänderungen nur mit hohem technischem Aufwand oder gar nicht praktikabel, sodass flexible, einfache Regelstrategien wie bei Asynchronmaschinen, die wir in späteren Kapitel betrachten werden, schwer umzusetzen sind.}
	\silence{2}
\end{frame}

\s{\subsection{Aufbau}\index{Synchronmaschine>Aufbau}}
\begin{frame}
	\ftx{Aufbau einer fremderregten Synchronmaschine}
	\begin{columns}
		\column[c]{0.55\textwidth}
		\begin{itemize}
			\item Drei Statorwicklungen, die 120° zueinander versetzt sind.
			\item Geblechte Ausführung des Stators, da dieser ein magnetisches Wechselfeld erfährt.
			\item Im Läufer befindet sich die Erregerwicklung.
			\item Der Läufer wird teilweise nicht geblecht.
		\end{itemize}
		\column[c]{0.45\textwidth}
		\s{
			\fo{width=0.7\textwidth, alt={Schematische Schnittansicht einer fremderregten Synchronmaschine. Das Bild zeigt die verschiedenen Komponenten des Motors, einschließlich des Gehäuses, der Lager, des Läufers, der Erregerwicklung sowie des Statorblechpakets und der Statorwicklungen.}}{}{Aufbau_Synchronmaschine1}{
				\put(83,0){\line(-1,1){37.5}}\put(83.3,0){Statorwicklungen}
				\put(83,6){\line(-1,1){33.5}}\put(83.3,6){Statorblechpaket}
				\put(89,18){\line(-1,1){32}}\put(89.3,18){Erregerwicklung}
				\put(89,24){\line(-1,1){35}}\put(89.3,24){Läufer}
				\put(100,36){\line(-1,1){26.2}}\put(100.3,36){Lager}
				\put(100,42){\line(-1,1){26.2}}\put(100.3,42){Gehäuse}
			}{{\bf Grafik einer aufgeschnittenen fremderregten Synchronmaschine.}}
			\vspace{0.5cm}
		}
		\b{
			\f{}{width=\textwidth}{Aufbau_Synchronmaschine1}{}
		}
	\end{columns}
	\speech{SynchronmaschineAufbau1}{1}{Eine Synchronmaschine besteht aus einem Stator mit drei symmetrisch angeordneten Wicklungen, die jeweils um 120 Grad zueinander versetzt sind und so ein drehendes magnetisches Feld erzeugen. Der Stator ist gegblecht ausgeführt, da er einem magnetischen Wechselfeld ausgesetzt ist und die Blechung Wirbelstromverluste reduziert. Im Läufer befindet sich die Erregerwicklung, die ein konstantes Feld erzeugt. Läufer sind dabei teils nicht geblecht ausgeführt, da hier kein wechselndes Magnetfeld in gleicher Stärke auftritt.}
	\silence{2}
\end{frame}

\begin{frame} \ftx{Aufbau einer fremderregten Synchronmaschine}
	\s{
		\fo{width=0.8\textwidth, alt={Querschnitt einer remderregten Synchronmaschine. In der Mitte befindet sich der Rotor mit der Erregerwicklung. Dieser wird ringförmig umschlossen vom Stator, der drei Statorwicklungen enthält. Diese sind jeweils um 120 Grad zueinander versetzt angeordnet.}}{}{Aufbau_Synchronmaschine2}{
			\put(32.3,16.1){\line(1,1){15}}\put(11,16.1){Statorwicklung 1}
			\put(86,66){\line(-1,-1){12.7}}\put(86.3,66){Statorwicklung 2}
			\put(91,10.4){\line(-1,1){12.4}}\put(91.3,10.4){Statorwicklung 3}
		}{{\bf Aufbau einer fremderregten Synchronmaschine.} Die um 120° zueinander versetzten Statorwicklungen sind hier farbig veranschaulicht.}
	}
	\b{
		\foo{}{width=0.75\textwidth}{Aufbau_Synchronmaschine2}{
			\put(32.3,16.1){\line(1,1){15}}\put(6.3,16.1){Statorwicklung 1}
			\put(86,66){\line(-1,-1){12.7}}\put(86.3,66){Statorwicklung 2}
			\put(91,10.4){\line(-1,1){12.4}}\put(91.3,10.4){Statorwicklung 3}
		}{}
	}
	\speech{SynchronmaschineAufbau2}{1}{Hier sehen wir den Querschnitt einer fremderregten Synchronmaschine. Es sind die drei Statorwicklungen erkennbar. Diese Wicklungen sind gleichmäßig um den Stator angeordnet und erzeugen durch Wechselstrom ein rotierendes Magnetfeld. Der Rotor, der durch die fremderregte Spannung magnetisiert wird, dreht sich synchron zur Frequenz des erzeugten Magnetfeldes im Stator. Diese Anordnung ermöglicht eine effiziente Umwandlung von elektrischer in mechanische Energie.}
	\silence{2}
\end{frame}



\begin{frame}\ftb{Feldverläufe}
	\s{
		\fo{width=\textwidth, alt={Das Diagramm zeigt drei sinusförmige Wellen, die die Ströme I1, I2 und I3 darstellen. Die Sinuskurven sind um 120 Grad zueinander versetzt. Unterhalb des Diagramms befindet sich jeweils an den Positionen 0, 120, 240 und 360 Grad eine Abbildung des Querschnitts einer fremderregten Synchronmaschine. Die in den Abbildungen eingezeichneten Flussrichtungen der Ströme innerhalb der drei Statorwickungen verhalten sich analog zu den im Diagramm dargestellten Sinuskurven.}}{}{Synchronmaschine_Feldverlaeufe}{
			\put(10.23,23.3){\begin{tikzpicture}[domain=0:36, samples=360, scale=0.982]
					\foreach \n in {0,...,3}{
							\draw[very thin, color=gray] (\n*4, 1) -- (\n*4,-1.5);
						}
					\draw (-0.1,0) -- (12.1,0);
					\draw[->] (0,-1.1) -- (0,1.2) node[above] {$I$};
					\draw (0.1, 0) node[below] {0};
					\foreach \n in {1,...,6}{
							\draw (\n*2, 0.1) -- (\n*2,-0.1);
							\fill[fill=white] (\n*2 - 0.1,-0.5) rectangle ++(0.2,0.4); % weißes Overlay über Linie hinter Zahlen
							\draw (\n*2, 0) node[below] {\pgfmathparse{60*\n}\pgfmathprintnumber[]{\pgfmathresult}\degree};
						}
					\draw[color=statorw1, scale=1/3, smooth, very thick] plot[id=drehfeldverlauf1] function{3 * cos(0.0556 * pi * x)};
					\draw[color=statorw1] (1.8,0.55) node[above] {$I_1$};
					\draw[color=statorw2, scale=1/3, smooth, very thick] plot[id=drehfeldverlauf3] function{3 * cos(0.0556 * pi * x - pi * 2/3)};
					\draw[color=statorw2] (5.8,0.55) node[above] {$I_2$};
					\draw[color=statorw3, scale=1/3, smooth, very thick] plot[id=drehfeldverlauf2] function{3 * cos(0.0556 * pi * x + pi * 2/3)};
					\draw[color=statorw3] (9.8,0.55) node[above] {$I_3$};
				\end{tikzpicture}
			}
			
			\put(12.5, 21.5){\color{statorw1}$I_1$} \put(38, 21.5){\color{statorw1}$I_1$} \put(63.6, 21.5){\color{statorw1}$I_1$} \put(89.3, 21.5){\color{statorw1}$I_1$}
			\put(21, 1.5){\color{statorw2}$I_2$} \put(47, 1.5){\color{statorw2}$I_2$} \put(73, 1.5){\color{statorw2}$I_2$} \put(98, 1.5){\color{statorw2}$I_2$}
			\put(0, 1.5){\color{statorw3}$I_3$} \put(25.5, 1.5){\color{statorw3}$I_3$} \put(51, 1.5){\color{statorw3}$I_3$} \put(77, 1.5){\color{statorw3}$I_3$}
		}{{\bf Richtung des Drehfeldvektors.} Der Rotor dreht sich synchron zu den sinusförmigen, um 120$\degree$ zueinander versetzten Drehfeldern des Stators.}
	}
	\b{
		\foo{}{height=0.8\textheight}{Synchronmaschine_Feldverlaeufe_a1}{
			\put(-47,60){\begin{tikzpicture}[domain=0:0.3, samples=360, scale=0.838]
					\foreach \n in {1,...,4}{
							\draw[very thin, color=gray] (\n*4, 1) -- (\n*4,-1);
						}
					\draw (-0.1,0) -- (12.1,0);
					\draw[->] (0,-1.1) -- (0,1.2) node[above] {$I$};
					\draw (0.2, 0) node[below] {0};
					\foreach \n in {1,...,6}{
							\draw (\n*2, 0.1) -- (\n*2,-0.1);
							\fill[fill=white] (\n*2 - 0.1,-0.6) rectangle ++(0.2,0.5); % weißes Overlay über Linie hinter Zahlen
							\draw (\n*2, 0) node[below] {\pgfmathparse{60*\n}\pgfmathprintnumber[]{\pgfmathresult}\degree};
						}
					\draw[color=statorw1, scale=1/3, smooth, very thick] plot[id=drehfeldverlauf1a1] function{3 * cos(0.0556 * pi * x)};
					\draw[color=statorw1] (-0.3,0.55) node[above] {$I_1$};
					\draw[color=statorw2, scale=1/3, smooth, very thick] plot[id=drehfeldverlauf3a1] function{3 * cos(0.0556 * pi * x - pi * 2/3)};
					\draw[color=statorw2] (-0.3,-0.1) node[above] {$I_2$};
					\draw[color=statorw3, scale=1/3, smooth, very thick] plot[id=drehfeldverlauf2a1] function{3 * cos(0.0556 * pi * x + pi * 2/3)};
					\draw[color=statorw3] (-0.3,-0.75) node[above] {$I_3$};
				\end{tikzpicture}
			}
			\put(28, 51){\color{statorw1}$I_1$}
			\put(48, 3){\color{statorw2}$I_2$}
			\put(-1, 3){\color{statorw3}$I_3$}
		}{}
	}
	\speech{SynchronmaschineFeldverlaufe1}{1}{Sehen wir uns das einmal Schritt für Schritt genauer an. Im Folgenden wird die Drehung des Rotors einer Synchronmaschine um 360 Grad in vier Schritten gezeigt. Die zeitlichen Verläufe der Ströme I-1, I-2 und I-3, die die Stromflüsse durch die drei Statorwicklungen darstellen, bewegen sich dabei sinusförmig. Im ersten Schritt, bei einer Rotordrehung von Null Grad, ist der Strom I-1 an seinem positiven Maximum, während die Stöme I-2 und I-3 negativ bei mittlerer Stärke ausfallen. Im Querschnitt des Rotos sind die Stromflüsse jweils durch Kreuze, für in die Bildfläche hinein bewegenden, und durch Punkte, für aus der Bildfläche hinaus bewegenden Strömen, dargestellt. Außerdem ist hier auch die Flussrichtung der jeweiligen Ströme an den Pfeilen abzulesen. Das von den drei Statorwicklungen erzeugte Magnetfeld, das letztlich den Rotor in Bewegung versetzt, wird mit orangefarbenen Linien angedeutet.}
	\silence{2}
\end{frame}

\b{ % Quasi-Animation für Vorlesungsdatei
	\begin{frame}\ftb{Feldverläufe}
		\foo{}{height=0.8\textheight}{Synchronmaschine_Feldverlaeufe_a2}{
			\put(-47,60){\begin{tikzpicture}[domain=0:12, samples=360, scale=0.838]
					\foreach \n in {1,...,4}{
							\draw[very thin, color=gray] (\n*4, 1) -- (\n*4,-1);
						}
					\draw (-0.1,0) -- (12.1,0);
					\draw[->] (0,-1.1) -- (0,1.2) node[above] {$I$};
					\draw (0.2, 0) node[below] {0};
					\foreach \n in {1,...,6}{
							\draw (\n*2, 0.1) -- (\n*2,-0.1);
							\fill[fill=white] (\n*2 - 0.1,-0.6) rectangle ++(0.2,0.5); % weißes Overlay über Linie hinter Zahlen
							\draw (\n*2, 0) node[below] {\pgfmathparse{60*\n}\pgfmathprintnumber[]{\pgfmathresult}\degree};
						}
					\draw[color=statorw1, scale=1/3, smooth, very thick] plot[id=drehfeldverlauf1a2] function{3 * cos(0.0556 * pi * x)};
					\draw[color=statorw1] (-0.3,0.55) node[above] {$I_1$};
					\draw[color=statorw2, scale=1/3, smooth, very thick] plot[id=drehfeldverlauf3a2] function{3 * cos(0.0556 * pi * x - pi * 2/3)};
					\draw[color=statorw2] (-0.3,-0.1) node[above] {$I_2$};
					\draw[color=statorw3, scale=1/3, smooth, very thick] plot[id=drehfeldverlauf2a2] function{3 * cos(0.0556 * pi * x + pi * 2/3)};
					\draw[color=statorw3] (-0.3,-0.75) node[above] {$I_3$};
				\end{tikzpicture}
			}
			\put(28, 51){\color{statorw1}$I_1$}
			\put(48, 3){\color{statorw2}$I_2$}
			\put(-1, 3){\color{statorw3}$I_3$}
		}{}
		\speech{SynchronmaschineFeldverlaufe2}{1}{Im nächsten Schritt sieht es nun anders aus: Wegen des sinusförmigen, um 120 Grad zueinander versetzten Stromverlaufs, haben sich die Stomstärken und Flussrichtungen verändert. Der Stom I-2 hat sein positives Maximum erreicht, während sich I-1 und I-2 bei mittlerer Stärke im negativen Bereich befinden. Diese Stomflüsse können wir auch im Querschnitt ablsenen, genauso wie das daraus resultierende Magnetfeld, das den Rotor um 120 Grad im Uhrzeigersinn weiter gedreht hat.}
	\end{frame}
	\begin{frame}\ftb{Feldverläufe}
		\foo{}{height=0.8\textheight}{Synchronmaschine_Feldverlaeufe_a3}{
			\put(-47,60){\begin{tikzpicture}[domain=0:24, samples=360, scale=0.838]
					\foreach \n in {1,...,4}{
							\draw[very thin, color=gray] (\n*4, 1) -- (\n*4,-1);
						}
					\draw (-0.1,0) -- (12.1,0);
					\draw[->] (0,-1.1) -- (0,1.2) node[above] {$I$};
					\draw (0.2, 0) node[below] {0};
					\foreach \n in {1,...,6}{
							\draw (\n*2, 0.1) -- (\n*2,-0.1);
							\fill[fill=white] (\n*2 - 0.1,-0.6) rectangle ++(0.2,0.5); % weißes Overlay über Linie hinter Zahlen
							\draw (\n*2, 0) node[below] {\pgfmathparse{60*\n}\pgfmathprintnumber[]{\pgfmathresult}\degree};
						}
					\draw[color=statorw1, scale=1/3, smooth, very thick] plot[id=drehfeldverlauf1a3] function{3 * cos(0.0556 * pi * x)};
					\draw[color=statorw1] (-0.3,0.55) node[above] {$I_1$};
					\draw[color=statorw2, scale=1/3, smooth, very thick] plot[id=drehfeldverlauf3a3] function{3 * cos(0.0556 * pi * x - pi * 2/3)};
					\draw[color=statorw2] (-0.3,-0.1) node[above] {$I_2$};
					\draw[color=statorw3, scale=1/3, smooth, very thick] plot[id=drehfeldverlauf2a3] function{3 * cos(0.0556 * pi * x + pi * 2/3)};
					\draw[color=statorw3] (-0.3,-0.75) node[above] {$I_3$};
				\end{tikzpicture}
			}
			\put(28, 51){\color{statorw1}$I_1$}
			\put(48, 3){\color{statorw2}$I_2$}
			\put(-1, 3){\color{statorw3}$I_3$}
		}{}
		\speech{SynchronmaschineFeldverlaufe3}{1}{Weitere 120 Grad weiter, bei einer Rotordrehung von 240 Grad, hat nun der Strom I-3 sein positives Maximum erreicht, während die anderen beiden Ströme bei mittlerer Stärke im negativen Bereich liegen.}
	\end{frame}
	\begin{frame}\ftb{Feldverläufe}
		\foo{}{height=0.8\textheight}{Synchronmaschine_Feldverlaeufe_a1}{
			\put(-47,60){\begin{tikzpicture}[domain=0:36, samples=360, scale=0.838]
					\foreach \n in {1,...,4}{
							\draw[very thin, color=gray] (\n*4, 1) -- (\n*4,-1);
						}
					\draw (-0.1,0) -- (12.1,0);
					\draw[->] (0,-1.1) -- (0,1.2) node[above] {$I$};
					\draw (0.2, 0) node[below] {0};
					\foreach \n in {1,...,6}{
							\draw (\n*2, 0.1) -- (\n*2,-0.1);
							\fill[fill=white] (\n*2 - 0.1,-0.6) rectangle ++(0.2,0.5); % weißes Overlay über Linie hinter Zahlen
							\draw (\n*2, 0) node[below] {\pgfmathparse{60*\n}\pgfmathprintnumber[]{\pgfmathresult}\degree};
						}
					\draw[color=statorw1, scale=1/3, smooth, very thick] plot[id=drehfeldverlauf1] function{3 * cos(0.0556 * pi * x)};
					\draw[color=statorw1] (-0.3,0.55) node[above] {$I_1$};
					\draw[color=statorw2, scale=1/3, smooth, very thick] plot[id=drehfeldverlauf3] function{3 * cos(0.0556 * pi * x - pi * 2/3)};
					\draw[color=statorw2] (-0.3,-0.1) node[above] {$I_2$};
					\draw[color=statorw3, scale=1/3, smooth, very thick] plot[id=drehfeldverlauf2] function{3 * cos(0.0556 * pi * x + pi * 2/3)};
					\draw[color=statorw3] (-0.3,-0.75) node[above] {$I_3$};
				\end{tikzpicture}
			}
			\put(28, 51){\color{statorw1}$I_1$}
			\put(48, 3){\color{statorw2}$I_2$}
			\put(-1, 3){\color{statorw3}$I_3$}
		}{}
		\speech{SynchronmaschineFeldverlaufe4}{1}{Bewegt sich die um sinusförmige Modulation der Ströme in den drei Erregerspulen um weitere 120 Grad, ist wie Ausgangsposition wieder erreicht und der Rotor hat eine 360 Grad Drehung im Uhrzeigersinn vollzogen. Diese Drehung wurde zur Veranschaulichung in vier Schritten dargestellt, da sich damit die Stromänderungen in den drei Statorwicklungen zu zeigen lassen. Tatsächlich müsste zwischen diesen Schritten jeweils eine quasi unendliche Zahl an Zwischenschritten gedacht werden, um die sinusförmigen Verläufe darzustellen.}
		\silence{2}
	\end{frame}
}


\begin{frame}\ftx{Drehmoment}
	\s{Bei der Synchronmaschine bewirkt die magnetische Kopplung zwischen Stator und Rotor das Drehmoment. Im Motorbetrieb folgt das Rotorfeld dem Statorfeld. Die mechanische Drehzahl muss daher immer gleich der elektrischen Drehfrequenz sein. Zwischen dem Stator und dem Rotor befindet sich der lastabhängige sogenannte Polradwinkel\index{Polradwinkel} $\vartheta$. \nomenclature[F]{$\vartheta$}{Polradwinkel der Synchronmaschine \nomunit{$\degree$}} Bei $\vartheta=90\degree$ ist das Drehmoment am größten. Wird im Motorbetrieb der Polradwinkel durch eine zu hohe Belastung größer als $90\degree$, sinkt das erzeugte Drehmoment und die Maschine \glqq kippt\grqq -- das heißt, sie bleibt schlagartig stehen.}
	\begin{itemize}
		\item Der Statorfluss entsteht durch die Überlagerung der Flüsse der drei Wicklungen.
		\item Der Läuferfluss wird durch die Erregerwicklung erzeugt.
	\end{itemize}
	\fu{
		\input{Tikz/src/drehmomentkennlinie_synchron}\alttext{Ein Diagramm, das zwei Kurven zeigt. Auf der X-Achse wird der Polradwinkel von 0 bis 180 Grad dargestellt. Auf der Y-Achse wird der Drehmoment abgebildet. Zwei parabelförmige Kurven UN und einhalb UN von 0 bix 180 Grad, deren Maximum bei 90 Grad liegt, sind in das Diargramm eingezeichnet. Das Maximum der Kurve von einhalb UN liegt etwa auf halber Höhe des Maximums von UN.}
	}{{\bf Drehmomentkennlinie in Abhängigkeit vom Polradwinkel.}}
	\speech{SynchronmaschineDrehmoment}{1}{Bei der Synchronmaschine bewirkt die magnetische Kopplung zwischen Stator und Rotor das Drehmoment. Im Motorbetrieb folgt das Rotorfeld dem Statorfeld. Die mechanische Drehzahl muss daher immer gleich der elektrischen Drehfrequenz sein. Zwischen dem Stator und dem Rotor befindet sich der lastabhängige sogenannte Polradwinkel Theta. Bei Theta gleich 90 Grad ist das Drehmoment am größten. Wird im Motorbetrieb der Polradwinkel durch eine zu hohe Belastung größer als 90 Grad, sinkt das erzeugte Drehmoment und die Maschine kippt, das heißt, sie bleibt schlagartig stehen.}
	\silence{2}
\end{frame}


\s{\subsection{Ersatzschaltbild eines Vollpolgenerators \index{Synchronmaschine>Ersatzschaltbild}}}
\begin{frame}
	\ftx{Einphasiges ESB Vollpolgenerator}
	\s{Im Netzbetrieb sind die Klemmenspannung $U_1$ und die Frequenz $\omega$ fest vorgegeben. Das ist ein häufiger Betriebsfall für die Synchronmaschine.}
	\b{Komplexes einphasiges ESB (Ersatzschaltbild) des Vollpolgenerators (z.B. im Netzbetrieb):}
	\fu{
		\input{Tikz/src/ESB_vollpolgenerator}\alttext{Schaltbild eines einphasigen Vollpolgenerators. Auf der linken Seite bedindet sich eine Reihenschaltung bestehend aus einem Widerstand R1, einer Induktivität X1 Sigma, einer Induktivität XH und einer Spannungsquelle Up, welche über die Spannungsquelle U1 geschlossen ist. Parallel zu U1 und Up ist zwischen X1 Sigma und Xh die Spannung Uq eingezeichntet. An der Anordnung liegt der Strom I1 an. Auf der rechten Seite bedindet sich eine Induktivität mit dem Strom IE.}
	}{{\bf Komplexes einphasiges ESB eines Vollpolgenerators.} Beispielsweise im Netzbetrieb.}
	\speech{SynchronmaschineVollpolgenerator}{1}{Wird eine Synchronmaschine nicht als Antrieb, sondern umgekehrt als Generator genutzt, spricht man vom Vollpolgenerator. Diese Generatoren zeichnen sich durch die Verwendung von Vollpolwicklungen aus, die eine gleichmäßige und symmetrische elektrische Leistung erzeugen. Bei Vollpolgeneratoren sind die Wicklungen so gestaltet, dass sie einen vollen magnetischen Pol ausnutzen, was bedeutet, dass jeweils die gesamte Poleinheit für die Stromerzeugung genutzt wird. Diese Anordnung führt zu einer optimalen Nutzung des magnetischen Flusses. Hier dargestellt ist ein Ersatzschaltbild eines einphasigen Vollpolgenerators. Im Netzbetrieb sind die Klemmenspannung U-1 und die Frequenz Omega fest vorgegeben. Das ist ein häufiger Betriebsfall für die Synchronmaschine.}
	\silence{2}
\end{frame}
